\chapter{Design and Methodology}\label{ch:B}

\section{System Architecture Overview}

The system consists of three loosely coupled components, each independently
deployable and testable:

\begin{enumerate}
  \item \textbf{Data pipeline} --- dataset curation, physics-based augmentation,
        feature engineering, and model training
        (\texttt{expand\_dataset.py}, \texttt{train.py},
        \texttt{utils/preprocessing.py}).

  \item \textbf{Backend API} --- a Flask REST service that loads trained model
        artifacts at startup and serves predictions with confidence scores
        (\texttt{app.py}).

  \item \textbf{Frontend} --- a single-page web application using the React
        framework and Vite for compiling and bundling the code, served
        statically via Nginx which also handles reverse proxying for the API.
\end{enumerate}

The backend was deployed as a WSGI server using Gunicorn, with two worker
processes running on an Ubuntu~24.04 machine.
Nginx sat in front and did two things --- it delivered the compiled React
bundle statically, and it forwarded all requests starting with \texttt{/api/}
back upstream to Gunicorn on port 5001.
To the client's browser, everything came from the same origin --- there was
never a cross-origin request and no port number to configure.

\section{Dataset}

\subsection{Collection and Structure}

Most of the time spent on assembling the dataset was during the first semester.
Using search terms such as \emph{``photoelectrochemical water splitting''},
\emph{``photoanode photocurrent density''}, and \emph{``STH efficiency semiconductor''},
we searched both Google Scholar and Web of Science.
All results had to be reviewed individually --- many studies utilised non-standard
illumination methods, provided only qualitative findings, or examined photocatalytic
suspensions instead of electrode-based systems --- all of which were excluded from
the dataset.
Of the 529 papers remaining after manual review, they represent approximately
twenty years of PEC research --- from early studies involving titanium dioxide
(TiO\textsubscript{2}), through to current high-performance bismuth vanadate
(BiVO\textsubscript{4}) and gallium nitride (GaN) materials.
Each row within the dataset represents a single experimental setup, identified by
DOI, taken from each individual study.

The thirty columns present information regarding the materials investigated,
structural properties, electrochemical test conditions, and various performance
metrics.
The most significant columns related to the investigation and their corresponding
data types are listed in Table~\ref{tab:columns}.

\begin{table}[htb]
\centering
\caption{Key dataset columns, types, and coverage in the original (529-row) dataset.}
\label{tab:columns}
\begin{tabular}{llr}
\toprule
Column & Type & Coverage (\%) \\
\midrule
Material System              & Categorical & 97 \\
Bandgap (eV)                 & Numerical   & 49 \\
Synthesis Method             & Categorical & 80 \\
Nanostructure                & Categorical & 67 \\
Cocatalyst                   & Categorical & 62 \\
Electrolyte                  & Categorical & 73 \\
pH                           & Numerical   & 47 \\
Light Intensity (mW/cm\textsuperscript{2}) & Numerical & 61 \\
Applied Bias (V vs.\ RHE)   & Numerical   & 74 \\
Temperature (K)              & Numerical   & 21 \\
Photocurrent Density         & Numerical   & 78 \\
STH Efficiency (\%)          & Numerical   & \textbf{9} \\
H\textsubscript{2} Evolution Rate & Numerical & \textbf{11} \\
\bottomrule
\end{tabular}
\end{table}

The values of row coverage correspond to the difficulty of the measurements.
In almost all papers, information about material composition and method of
preparation can be found --- it is a standard element of the experimental
description.
Photocurrent density is mentioned in 78\% of cases, since it is a standard output
of practically any PEC experiment.
STH and H\textsubscript{2} coverage then fall to 9\% and 11\%.
Measuring STH requires knowledge of the incident radiation power and Faradaic
efficiency.
Determining H\textsubscript{2} evolution rate requires a gas-tight electrochemical
cell, which most laboratories publishing photocurrent data cannot afford to use on
a routine basis.
Therefore the lack of appropriate equipment leads to an instrumentation gap, which
causes the data scarcity addressed by the augmentation pipeline developed in this
work.
Only 48 STH rows and 57 H\textsubscript{2} rows exist; therefore even a 5-fold
cross-validation results in as few as 10 validation samples per fold for estimating
$R^2$ values.

\subsection{Physics-Based Dataset Augmentation}

To address the sparsity problem, two complementary augmentation strategies were
applied, producing a final dataset of \textbf{676 rows}.

\subsubsection{Strategy 1: Physics-Derived Estimates for Existing Rows}

For the 364 rows in which photocurrent density has been determined via measurement
but STH has not been reported, the STH value was estimated using the standard
approximation (Equation~\eqref{eq:sth}):

\begin{equation}
  \widehat{\mathrm{STH}} =
    \frac{J_{ph} \times 1.23}{P_{in}} \times \eta_F \times 100\%
  \label{eq:sth_est}
\end{equation}

where $P_{in}$ is set to the reported light intensity when available, or to
100\,mW/cm\textsuperscript{2} (1\,sun, AM~1.5G standard) otherwise;
$\eta_F$ is sampled uniformly from $[0.75,\,0.98]$ for unbiased measurements
or $[0.55,\,0.85]$ for biased measurements ($|V_{bias}| > 0.1$\,V) to reflect
realistic Faradaic losses.
A multiplicative noise factor $\delta \sim \mathcal{U}(0.92,\,1.08)$ is applied
to reproduce typical measurement variability across different laboratories.

Similarly, the H\textsubscript{2} evolution rate was derived from Faraday's law
(Equation~\eqref{eq:her}) with $\eta_F \sim \mathcal{U}(0.60,\,0.95)$.

This filled \textbf{361} STH values and \textbf{357} H\textsubscript{2} values.
The scientific justification is that photocurrent and STH are not independent
quantities --- STH is defined as a function of photocurrent via a fixed
thermodynamic formula.
A paper that reports $J_{ph} = 3.5$\,mA/cm\textsuperscript{2} under standard
illumination has implicitly already determined $\mathrm{STH} \approx 4.3\%$,
whether or not the authors computed it.
We are recovering a value already encoded in the data, not fabricating a new
measurement.
The Faradaic noise term prevents the recovered labels from becoming a perfectly
deterministic function of photocurrent, which would otherwise allow the STH model
to learn a trivial linear mapping rather than the material-dependent relationships
we care about.

\subsubsection{Strategy 2: Literature-Based Synthetic Rows}

The second strategy added 147 entirely new rows for 25 well-studied material systems.
For each material, photocurrent ranges were taken from published review
articles~\cite{dias2016recent,tamirat2016using} --- these are ranges that have
been measured and reported many times across independent labs, so they represent
a reasonable envelope of achievable performance.
Structural parameters and operating conditions were sampled to match typical
experimental practice for each material: BiVO\textsubscript{4}, for example,
was paired with near-neutral Na\textsubscript{2}SO\textsubscript{4} electrolyte
at pH~6--7 and a modest applied bias, consistent with how it is almost always
tested in the literature.
The covered systems include BiVO\textsubscript{4} variants,
$\alpha$-Fe\textsubscript{2}O\textsubscript{3}, WO\textsubscript{3},
TiO\textsubscript{2}, ZnO, Cu\textsubscript{2}O, GaN,
Ta\textsubscript{3}N\textsubscript{5}, Si, g-C\textsubscript{3}N\textsubscript{4},
CuInS\textsubscript{2}, and several ternary oxides.
Every synthetic row carries the tag \texttt{Data\_Source = "synthetic"} so it
can be identified and excluded from any analysis that requires only real measurements.

\subsection{Missing Value Imputation}

Before running augmentation, four supporting fields were filled using deterministic
rules rather than statistical imputation:

\begin{itemize}
  \item \textbf{Bandgap}: for 204 rows with known material but missing bandgap,
        reference values were drawn from a lookup table of 40+ materials
        (e.g., TiO\textsubscript{2} anatase: 3.2\,eV; BiVO\textsubscript{4}: 2.4\,eV;
        $\alpha$-Fe\textsubscript{2}O\textsubscript{3}: 2.1\,eV) with a small jitter
        $\delta \sim \mathcal{U}(-0.05,\,+0.05)$\,eV to prevent the lookup
        from creating artificial clusters.

  \item \textbf{pH}: inferred from electrolyte composition for 114 rows
        (KOH / NaOH $\to$ pH~13--14;
         H\textsubscript{2}SO\textsubscript{4} $\to$ pH~0--1.5;
         Na\textsubscript{2}SO\textsubscript{4} $\to$ pH~5.8--7.2).

  \item \textbf{Temperature}: set to 298\,K for all 420 rows lacking this value,
        consistent with standard laboratory ambient conditions.

  \item \textbf{Light intensity}: set to 100\,mW/cm\textsuperscript{2} for
        61 rows that specify AM~1.5G illumination but omit the numeric value.
\end{itemize}

\subsection{Final Dataset Statistics}

After augmentation, the dataset contains 676 rows.
Table~\ref{tab:augmented} compares key coverage metrics before and after
augmentation.

\begin{table}[htb]
\centering
\caption{Dataset coverage before and after physics-based augmentation.}
\label{tab:augmented}
\begin{tabular}{lrrrr}
\toprule
Target / Feature & Before & Before (\%) & After & After (\%) \\
\midrule
Photocurrent Density & 410 & 78 & 557 & 82 \\
STH Efficiency       & 48  &  9 & 556 & 82 \\
H\textsubscript{2} Evolution Rate & 57 & 11 & 559 & 83 \\
Bandgap (eV)         & 258 & 49 & 609 & 90 \\
pH                   & 250 & 47 & 509 & 75 \\
Temperature (K)      & 112 & 21 & 676 & 100 \\
\bottomrule
\end{tabular}
\end{table}

\section{Feature Engineering}

Twenty numeric features are derived from eight raw input fields.
Table~\ref{tab:features} lists each feature together with its physical
motivation.
The interaction features in the lower half of the table encode the most
important cross-variable relationships identified in the PEC literature
(Section~\ref{ch:A}).

\begin{table}[htb]
\centering
\caption{Physics-informed engineered features used by the ML models.}
\label{tab:features}
\begin{tabular}{lp{8cm}}
\toprule
Feature & Physical motivation \\
\midrule
\texttt{photon\_energy}       & $1240/E_g$ --- absorption edge wavelength (nm) \\
\texttt{overpotential\_proxy} & $V_{bias} - (E_g - 1.23)$: net driving force beyond thermodynamic minimum \\
\texttt{bandgap\_sq}          & $E_g^2$: non-linear dependence of absorption on bandgap \\
\texttt{log\_bandgap}         & $\ln(1+E_g)$: compressed bandgap representation \\
\texttt{ph\_deviation}        & $|pH - 7|$: distance from neutral (surface charge state proxy) \\
\texttt{bias\_positive}       & $\max(V_{bias},\,0)$: forward-bias contribution only \\
\texttt{temp\_normalized}     & $(T - 298)/298$: normalised deviation from room temperature \\
\texttt{has\_cocatalyst}      & Binary: 1 if cocatalyst is present, 0 otherwise \\
\texttt{is\_nanostructured}   & Binary: 1 if nanostructure is non-bulk \\
\texttt{ionic\_strength\_proxy} & Estimated ionic strength from electrolyte identity \\
\texttt{bandgap\_x\_ph}       & $E_g \times \mathrm{pH}$: flat-band potential shift (Nernst: $-59$\,mV/pH) \\
\texttt{bias\_x\_ionic}       & $V_{bias} \times I_{proxy}$: electric field in the double layer \\
\texttt{bg\_x\_overpot}       & $E_g \times \eta_{proxy}$: combined thermodynamic driving force \\
\texttt{ph\_x\_temp}          & $\mathrm{pH} \times T$: Arrhenius-like surface kinetics term \\
\bottomrule
\end{tabular}
\end{table}

Six categorical features --- material system, electrolyte, synthesis method,
photoelectrode type, nanostructure, and morphology --- are one-hot encoded inside
a scikit-learn \texttt{Pipeline} with \texttt{handle\_unknown="ignore"} so that
unseen material names at inference time do not cause failures.
Median imputation is applied to all numeric features and mode imputation to all
categorical features within the same pipeline, ensuring consistent behaviour
between training and deployment.

\section{Data Leakage Prevention}\label{sec:leakage}

STH efficiency is defined as a deterministic function of photocurrent
(Equation~\eqref{eq:sth_est}).
Including photocurrent as a predictor for STH would inflate the cross-validated
$R^2$ to near 1.0 without the model learning any generalisable relationship.
The STH model therefore excludes photocurrent from its feature set
(\texttt{TARGET\_EXCLUSIONS["sth"] = ["photocurrent"]}).

Similarly, 33 photocathode entries --- which by measurement convention report
negative photocurrent values --- are removed from the photocurrent training set.
Retaining them would create a bimodal target distribution incompatible with
standard regression objectives.
These rows are preserved in the dataset for completeness but excluded from
the photocurrent target column during training.

\section{ML Model Pipeline}

\subsection{Model Selection Rationale}\label{sec:model_selection}

Two complementary tree ensemble methods are trained independently for each target.
As discussed in Section~\ref{ch:A}, tree-based ensembles are the empirically
superior choice for tabular scientific datasets of this scale~\cite{grinsztajn2022tree}.

\textbf{XGBoost Regressor} was configured with: 300 estimators, learning rate
$\eta = 0.05$, maximum tree depth 4, subsampling ratio 0.8 (rows) and
0.8 (features), $L_1$ regularisation $\alpha = 0.1$, $L_2$ regularisation
$\lambda = 1.0$, and minimum child weight 3.
The shallow trees (depth 4) and strong regularisation are deliberate: they
reduce the risk of overfitting to measurement artefacts in the heterogeneous
training data.

\textbf{Random Forest Regressor} was configured with: 300 estimators, maximum
depth 8, minimum samples per leaf 2, and the square-root feature selection
heuristic at each split.
The deeper trees (depth 8) without regularisation exploit a higher proportion
of the training signal at the cost of increased variance, which is controlled
by the averaging operation of the ensemble.

Both models are wrapped in scikit-learn \texttt{Pipeline} objects that include
the \texttt{ColumnTransformer} preprocessor described above.
The full pipeline is cloned and re-trained on all available data after
cross-validation evaluation to produce the final inference artifact.

\subsection{Ensemble Inference and Confidence Scoring}\label{sec:ensemble}

At prediction time, both saved models are queried independently.
The ensemble prediction is the CV-R\textsuperscript{2}-weighted average:

\begin{equation}
  \hat{y}_{ens} =
    \frac{\displaystyle\sum_{m} \max(0,\,r^2_m)\,\hat{y}_m}
         {\displaystyle\sum_{m} \max(0,\,r^2_m)}
  \label{eq:ensemble}
\end{equation}

This weighting scheme assigns near-zero influence to any model with a negative
or near-zero $R^2$, preventing a poor learner from corrupting the ensemble output.

Prediction uncertainty is characterised by the \emph{relative standard deviation}
of the two models' outputs:
\begin{equation}
  \sigma_{rel} = \frac{\mathrm{std}(\hat{y}_1,\hat{y}_2)}{\hat{y}_{ens}}
\end{equation}

The \emph{confidence score} combines model quality and inter-model agreement
into a single $[0,1]$ value:
\begin{equation}
  c = 0.6\,\bar{r}^2 + 0.4\,\max\!\left(0,\,1 - \sigma_{rel}\right)
  \label{eq:confidence}
\end{equation}
where $\bar{r}^2$ is the mean cross-validated $R^2$ clipped to $[0,1]$.
A score of 1.0 would require a perfect model ($r^2 = 1$) with zero disagreement
between the two estimators.
In practice, scores above 0.6 indicate that the prediction region is well covered
by training data and both models agree on the output magnitude.

\subsection{Performance Classification}\label{sec:classification}

To turn three continuous predictions into a single actionable label, we compute
a composite score on a 0--100 scale from the predicted photocurrent and STH values:
\begin{equation}
  s = \bigl(p_{pc} + b_{STH}\bigr)
      \times \bigl(0.4 + 0.6\,c\bigr)
  \label{eq:score}
\end{equation}
where:
\begin{itemize}
  \item $p_{pc} \in [0,\,100]$ is the photocurrent's percentile position
        relative to the training distribution: $p_{pc} = 100$ for
        $J_{ph} \geq p_{75}$, $p_{pc} = 50$ for $J_{ph} = p_{40}$,
        with linear interpolation between these thresholds.
  \item $b_{STH} \in \{0,\,7,\,15\}$ is a bonus for STH performance
        (15 points if STH $\geq p_{75}$, 7 if STH $\geq p_{40}$, 0 otherwise).
  \item $c$ is the confidence from Equation~\eqref{eq:confidence}.
\end{itemize}

The confidence factor $(0.4 + 0.6\,c)$ scales the score down for uncertain
predictions: a model with $r^2 = 0$ can contribute at most 40 points to the
composite score regardless of what the raw predicted values say, so low-confidence
predictions cannot reach the HIGH class.
The final thresholds --- HIGH ($s \geq 70$), MEDIUM ($40 \leq s < 70$),
LOW ($s < 40$) --- were chosen so that HIGH roughly corresponds to the top
quartile of the training distribution.

\subsection{Cross-Validation Protocol}

All models were evaluated with 5-fold cross-validation on the augmented dataset.
For each fold, training and validation splits were created without any data leakage ---
feature engineering was applied inside the pipeline, not beforehand.
The random seed was fixed at 42 for all splits and model initialisations so that
results can be reproduced exactly.
After CV scoring, each model was retrained on the full available data for that
target before saving, since the deployed model should use every row it can.

For the pre-augmentation results in Table~\ref{tab:cv_results}, the same procedure
was applied but with fold counts reduced for sparse targets.
At 48 STH rows, a strict 5-fold split leaves only about 9 samples per validation
fold --- too few for $R^2$ to be stable.
Reducing the fold count in this case gives a more honest estimate of
generalisation ability.

\section{System Implementation}

\subsection{Training Pipeline}

The training script (\texttt{train.py}) runs in four stages.
It starts by loading the CSV and standardising column names --- different versions
of the dataset used slightly different capitalisations and spacing, so this
normalisation step runs before anything else.

\begin{enumerate}
  \item Load the CSV and standardise raw column names to internal identifiers.
  \item Parse and clean all numeric fields: values recorded as
        ``$1.4 \pm 0.2$'' are converted to their central value, and physically
        impossible inputs (negative photocurrent, pH outside $[0,14]$,
        bandgap below 0.5\,eV) are flagged and removed.
  \item Run the feature engineering module to compute all 20 derived columns
        via \texttt{utils/preprocessing.py}.
  \item Train both XGBoost and Random Forest for each target, run
        cross-validation, select the better model, and save the artifact
        as a \texttt{joblib} pickle containing the fitted pipelines,
        CV metrics, training-set medians and modes for inference-time
        imputation, and percentile thresholds for performance classification.
\end{enumerate}

\subsection{REST API}

The backend has three endpoints.

\begin{itemize}
  \item \texttt{GET /} --- a health check that returns the names of the
        loaded models and their CV $R^2$ scores.
        The frontend fetches this on startup to display model quality
        information before the user inputs anything.

  \item \texttt{POST /predict} --- takes a JSON body with up to 14 fields,
        runs feature engineering and inference, and returns predictions for
        all three targets with per-target confidence scores and the composite
        performance classification.

  \item \texttt{GET /dataset} --- returns the full training dataset as a
        JSON array for the frontend's data explorer tab.
\end{itemize}

Every input is validated against physical bounds before inference runs ---
pH must be between 0 and 14, bandgap between 0.5 and 6\,eV,
bias between $-2$ and 5\,V vs.\ RHE.
A call to validate that will fail returns an HTTP~400 response with the exact
name of the offending field, so the model does not create nonsensical outputs
when it encounters invalid physical input parameters.

\subsection{Web Application}

The frontend application is built using React and was compiled into a statically
bundled version of the code by Vite, served in production on port~80 via Nginx.
All interactions made against the Flask API are proxied through Nginx, so users
have no direct interaction with the Flask service itself.
There are four pages within this interface.

\begin{itemize}
  \item \textbf{Dashboard} --- when opened, loads immediately and presents the
        same quick prediction form alongside the current model $R^2$ scores
        that are loaded dynamically each time the page is accessed.
        This allows the user to determine how well they can trust the output
        prior to entering any data.
  \item \textbf{Prediction} --- presents the complete parameter prediction form,
        placing the five required fields in the upper portion and the remainder
        (optional) at the lower end.
        This keeps the form functional even when it may appear overwhelming to
        a new user.
  \item \textbf{Results} --- displays the three values generated by prediction,
        their confidence percentage, and colour-coded performance bars
        (green for HIGH, yellow/orange for MEDIUM, red for LOW).
  \item \textbf{About} --- shows the feature importance table, cross-validation
        results, and composition breakdowns of the datasets used to build the
        system.
\end{itemize}

All categorical input fields use \texttt{<datalist>} elements to offer
autocomplete suggestions while still permitting free-text entry, so
researchers can input any material or electrolyte not present in the
preset list.
